\begin{abstract}
This study evaluated whether students' answers about generative artificial intelligence use improved prediction of current burnout scores beyond study conditions alone. The analysis used 97 responses from a cross-sectional survey targeting Lebanese university students. The burnout outcome was a 12-item Burnout Assessment Tool for Students score on a 1-to-5 scale. Ridge regression was evaluated using five study-condition questions alone and using those same questions plus three on use frequency, dependence on these tools to complete academic work, and reliance on them in place of independent thinking. With the same evaluation splits, adding the three questions reduced mean absolute error from 0.6183 to 0.5883 points. For the High group, defined by a measured score of at least 4, the ridge model with all eight inputs had an error of 0.9329 points. The warning task compared logistic regression using workload alone with a random forest using all eight inputs. At the selected warning thresholds, the models detected 16.0 and 14.8 of the 17 High students on average, respectively. Both averaged 28.0 false warnings among the 80 students below the cutoff. The forest achieved 76.0\% balanced accuracy with the original labels, higher than in all 199 reruns with shuffled labels. Both tasks used repeated nested cross-validation within the same survey sample. An independent cohort is needed to assess performance on new students.
\end{abstract}
